\documentclass[11pt,a4paper]{article}
% Common preamble for the hint sheets (\input from each file)
\usepackage[utf8]{inputenc}
\usepackage[T1]{fontenc}
\usepackage{lmodern}
\usepackage{amsmath,amssymb,amsthm,mathtools}
\usepackage[a4paper,margin=2.3cm]{geometry}
\usepackage{enumitem}
\usepackage{xcolor}
\usepackage[hidelinks]{hyperref}
\renewcommand{\leq}{\leqslant}
\renewcommand{\le}{\leqslant}
\renewcommand{\geq}{\geqslant}
\renewcommand{\ge}{\geqslant}
\setlength{\parindent}{0pt}
\setlength{\parskip}{0.35em}
% \hintstitle{Sheet title}{Class work / Homework}
\newcommand{\hintstitle}[2]{\begin{center}{\LARGE\bfseries Hints: #1\par}\vspace{0.5em}{\large #2\par}\end{center}
\noindent\textit{The hints to each problem are cumulative. Read one, try again, and open the next only if you are still stuck.}\par\medskip\hrule\medskip}
\newcommand{\problem}[1]{\par\bigskip{\large\bfseries Problem #1}\par\nopagebreak\smallskip}
\newcommand{\subpart}[1]{\par\smallskip\textbf{(#1)}\par\nopagebreak}
\newcommand{\hint}[1]{\par\smallskip\noindent\textcolor{black!70}{\textbf{Hint #1.}}\enspace}
\newcommand{\note}[1]{\par\smallskip\noindent\textit{Note.} #1\par}
\providecommand{\lt}{<}
\providecommand{\gt}{>}
\providecommand{\tightlist}{\setlength{\itemsep}{0pt}\setlength{\parskip}{0pt}}
\emergencystretch=3em

\begin{document}
\hintstitle{Integrals -- 1. Substitutions: rationalisation, Euler, the half-angle substitution}{Homework}

\problem{1}
\hint{1}
For \(\int\sqrt{\tan x}\,dx\) put \(t=\sqrt{\tan x}\), i.e.~\(x=\arctan t^2\), \(dx=\frac{2t\,dt}{1+t^4}\); you get \(\int\frac{2t^2}{1+t^4}dt\). For the second integral divide numerator and denominator by \(\cos^4x\): with \(t=\tan x\) it becomes \(\int\frac{(1+t^2)\,dt}{1+t^4}\).
\hint{2}
The trick for \(\frac{1\pm t^2}{1+t^4}\): divide by \(t^2\) to get \(\frac{1\pm 1/t^2}{t^2+1/t^2}\), and note \(t^2+\frac1{t^2}=(t-\frac1t)^2+2=(t+\frac1t)^2-2\). So \(u=t-\frac1t\) or \(u=t+\frac1t\) reduces everything to \(\int\frac{du}{u^2\pm 2}\), a quadratic. Write \(2t^2=(t^2+1)+(t^2-1)\) for the first integral.
\hint{3}
Answers: \(\int\frac{dx}{\sin^4x+\cos^4x}=\frac1{\sqrt2}\arctan\frac{\tan^2x-1}{\sqrt2\tan x}+C\) (equivalently \(\frac1{\sqrt2}\arctan\frac{\tan2x}{\sqrt2}\), up to constants on intervals). With \(t=\sqrt{\tan x}\): \(\int\sqrt{\tan x}\,dx=\frac1{\sqrt2}\arctan\frac{t-1/t}{\sqrt2}+\frac1{2\sqrt2}\ln\frac{t+1/t-\sqrt2}{t+1/t+\sqrt2}+C\). With \(s=\tan\frac x2\) instead, \(\sin^4x+\cos^4x=\frac{16s^4+(1-s^2)^4}{(1+s^2)^4}\), a quartic in \(s^2\) in the denominator: much worse.

\problem{2}
\hint{1}
If \(P(x)=\prod(x-\omega_k)\) has simple roots, then \(\frac1{P(x)}=\sum_k\frac{1}{P'(\omega_k)(x-\omega_k)}\). For \(P=x^n+1\): \(P'(\omega)=n\omega^{n-1}=-n/\omega\), so \(\frac1{x^n+1}=-\frac1n\sum_k\frac{\omega_k}{x-\omega_k}\); for \(x^6-1\) similarly \(\frac1{x^6-1}=\frac16\sum\frac{\zeta_k}{x-\zeta_k}\).
\hint{2}
Pair conjugate roots \(e^{\pm i\theta}\): \(\frac{e^{i\theta}}{x-e^{i\theta}}+\frac{e^{-i\theta}}{x-e^{-i\theta}}=\frac{2x\cos\theta-2}{x^2-2x\cos\theta+1}\), which integrates to a \(\log\) plus an \(\arctan\). This is the general answer for \(\int\frac{dx}{x^n+1}\), with \(\theta=\frac{(2k+1)\pi}{n}\) (plus the real root \(-1\) when \(n\) is odd).
\hint{3}
\(\int\frac{dx}{x^4+1}=\frac1{4\sqrt2}\ln\frac{x^2+\sqrt2x+1}{x^2-\sqrt2x+1}+\frac1{2\sqrt2}\left(\arctan(\sqrt2x+1)+\arctan(\sqrt2x-1)\right)+C\). For \(x^6-1\): \(\frac1{x^6-1}=\frac16\left(\frac1{x-1}-\frac1{x+1}+\frac{x-2}{x^2-x+1}-\frac{x+2}{x^2+x+1}\right)\), giving \(\frac16\ln\left\lvert\frac{x-1}{x+1}\right\rvert+\frac1{12}\ln\frac{x^2-x+1}{x^2+x+1}-\frac1{2\sqrt3}\left(\arctan\frac{2x-1}{\sqrt3}+\arctan\frac{2x+1}{\sqrt3}\right)+C\).

\problem{3}
\hint{1}
\begin{enumerate}
\def\labelenumi{(\alph{enumi})}
\tightlist
\item
  \(u=2+4x\); (b) multiply by \(\frac{\sec x+\tan x}{\sec x+\tan x}\); (c) the numerator is half the derivative of the denominator; (d) \(\frac1{(x+1/x)^2}=\frac{x^2}{(x^2+1)^2}\); (e) divide: \(\frac{x^3}{1+x^2}=x-\frac{x}{1+x^2}\); (f) factor numerator and denominator by grouping.
\end{enumerate}
\hint{2}
(f): \(x^5-x^3+x^2-1=(x^2-1)(x^3+1)\) and \(x^4-x^3+x-1=(x-1)(x^3+1)\), so the integrand is \(x+1\). (d): \(\frac{x^2}{(x^2+1)^2}=\frac1{x^2+1}-\frac1{(x^2+1)^2}\), or integrate by parts with \(u=x\).
\hint{3}
Answers (\(+C\)): (a) \(\frac1{40}(2+4x)^{5/2}-\frac1{12}(2+4x)^{3/2}\); (b) \(\ln\lvert\sec x+\tan x\rvert\); (c) \(\frac12\ln(x^2+2x+3)\); (d) \(\frac12\left(\arctan x-\frac{x}{x^2+1}\right)\); (e) \(\frac{x^2}2-\frac12\ln(1+x^2)\); (f) \(\frac{x^2}2+x\).

\problem{4}
\hint{1}
\(\sqrt{1+\tan^2x}=\sqrt{\sec^2x}=\frac1{\lvert\cos x\rvert}\), not \(\sec x\). So the integrand is \(\frac{\sin x}{\lvert\cos x\rvert}=\pm\tan x\) depending on the sign of \(\cos x\).
\hint{2}
On intervals where \(\cos x\gt 0\) the antiderivative is \(-\ln\lvert\cos x\rvert+C\); where \(\cos x\lt 0\) it is \(+\ln\lvert\cos x\rvert+C\). In one formula: \(-\operatorname{sgn}(\cos x)\ln\lvert\cos x\rvert+C\) on each interval. The naive answer \(-\ln\lvert\cos x\rvert\) silently used \(\sqrt{\sec^2x}=\sec x\), true only where \(\cos x\gt 0\).

\problem{5}
\hint{1}
For \(I_n\), integrate \(I_{n-1}=\int\frac{dx}{(x^2+a^2)^{n-1}}\) by parts with \(u=(x^2+a^2)^{-(n-1)}\), \(dv=dx\), and write \(x^2=(x^2+a^2)-a^2\) in the new integral. For \(J_n\), write \(\sin^nx=\sin^{n-1}x\cdot\sin x\), take \(dv=\sin x\,dx\), and use \(\cos^2=1-\sin^2\).
\hint{2}
You should get \(I_n=\frac{x}{2(n-1)a^2(x^2+a^2)^{n-1}}+\frac{2n-3}{2(n-1)a^2}I_{n-1}\) with \(I_1=\frac1a\arctan\frac xa\), and \(J_n=-\frac{\sin^{n-1}x\cos x}{n}+\frac{n-1}{n}J_{n-2}\) with \(J_0=x\).
\hint{3}
\(I_3=\frac{x}{4a^2(x^2+a^2)^2}+\frac{3x}{8a^4(x^2+a^2)}+\frac{3}{8a^5}\arctan\frac xa+C\) and \(J_6=-\frac{\sin^5x\cos x}{6}-\frac{5\sin^3x\cos x}{24}-\frac{5\sin x\cos x}{16}+\frac{5x}{16}+C\).

\problem{6}
\hint{1}
Make the inverse function (or the inner root) disappear by differentiating it: (b), (c), (d), (e) take \(u=\) the inverse function and \(dv=dx\) (or \(x\,dx\)). In (a) substitute \(t=\sqrt x\) first; in (f) substitute \(t=e^x\) first.
\hint{2}
\begin{enumerate}
\def\labelenumi{(\alph{enumi})}
\tightlist
\item
  becomes \(2\int t^2e^t\,dt\); (f) becomes \(\int t^2\arctan t\,dt\), then parts with \(u=\arctan t\) and divide \(\frac{t^3}{1+t^2}=t-\frac t{1+t^2}\). (c) needs parts twice: the second time with \(u=\arcsin x\), \(dv=\frac{x\,dx}{\sqrt{1-x^2}}\). In (d) the leftover is \(\int\frac{\sqrt x}{2(1+x)}dx\), done by \(x=t^2\).
\end{enumerate}
\hint{3}
Answers (\(+C\)): (a) \(2e^{\sqrt x}(x-2\sqrt x+2)\); (b) \(\frac{x^2+1}2\arctan x-\frac x2\); (c) \(x\arcsin^2x+2\sqrt{1-x^2}\arcsin x-2x\); (d) \((x+1)\arctan\sqrt x-\sqrt x\); (e) \(x\operatorname{arcosh}x-\sqrt{x^2-1}\); (f) \(\frac{e^{3x}}3\arctan e^x-\frac{e^{2x}}6+\frac16\ln(1+e^{2x})\).

\problem{7}
\hint{1}
Rule: if \(R(-\sin x,\cos x)=-R(\sin x,\cos x)\) use \(t=\cos x\); if \(R(\sin x,-\cos x)=-R\) use \(t=\sin x\); if \(R(-\sin x,-\cos x)=R\) use \(t=\tan x\). (Reason: in the first case \(R\,dx=R_1(\cos x,\sin^2x)\sin x\,dx\), etc.)
\hint{2}
\(\frac{\sin^3x}{1+\cos^2x}\) is odd in \(\sin\): \(t=\cos x\) gives \(\int\frac{t^2-1}{t^2+1}dt\), so \(\cos x-2\arctan(\cos x)+C\). \(\frac1{1+\cos^2x}\) is even: \(t=\tan x\) gives \(\int\frac{dt}{t^2+2}=\frac1{\sqrt2}\arctan\frac{\tan x}{\sqrt2}+C\). \(\frac{\cos x}{2-\cos^2x}=\frac{\cos x}{1+\sin^2x}\) is odd in \(\cos\): \(t=\sin x\) gives \(\arctan(\sin x)+C\).

\problem{8}
\hint{1}
Complete the square and use a trigonometric or hyperbolic substitution: \(x-x^2=\frac14-(x-\frac12)^2\); \(12-3x^2=3(4-x^2)\) with \(x=2\sin\theta\); \(\sqrt{x^2-1}\) with \(x=\sec\theta\) (or \(\cosh u\)); \(\sqrt{x^2+25}\) with \(x=5\sinh u\) (or \(5\tan\theta\)).
\hint{2}
Answers (\(+C\)): \(\arcsin(2x-1)\) (equivalently \(2\arcsin\sqrt x\)); \(\sqrt3\left(\frac x2\sqrt{4-x^2}+2\arcsin\frac x2\right)\); for \(x\gt 1\), \(\int\tan^2\theta\,d\theta\) gives \(\sqrt{x^2-1}-\operatorname{arcsec}x\); \(\frac x2\sqrt{x^2+25}+\frac{25}2\ln\left(x+\sqrt{x^2+25}\right)\).

\problem{9}
\hint{1}
By parts twice, \(I=\int e^{ax}\cos bx\,dx\) reappears on the right with coefficient \(-\frac{b^2}{a^2}\); solve the resulting linear equation for \(I\). Alternatively \(\int e^{(a+ib)x}dx=\frac{e^{(a+ib)x}}{a+ib}\); multiply by \(\frac{a-ib}{a^2+b^2}\) and take the real part.
\hint{2}
Both give \(\frac{e^{ax}(a\cos bx+b\sin bx)}{a^2+b^2}+C\), and the imaginary part gives \(\int e^{ax}\sin bx\,dx\) for free; the two integrations by parts are exactly the real and imaginary parts of the single complex computation.

\problem{10}
\hint{1}
A denominator \(v^2\) suggests \((u/v)'=\frac{u'v-uv'}{v^2}\); a fraction whose numerator is the derivative of its denominator integrates to \(\ln\lvert v\rvert\). In (a) and (e) try \(v=\log x\), \(u=\sin x\). In (b) divide numerator and denominator by \(x\).
\hint{2}
\begin{enumerate}
\def\labelenumi{(\alph{enumi})}
\setcounter{enumi}{1}
\tightlist
\item
  \(\frac{x+1}{x(x+\log x)}=\frac{1+1/x}{x+\log x}\); (c) \(1+\cos x=(x+\sin x)'\). (d) is not quite of this form: write \(1+2x^{2022}=(1+x^{2022})+x^{2022}\), so the integrand is \(\frac1x+\frac{x^{2021}}{1+x^{2022}}\). (e) equals \(\frac{\cos x}{\log x}-\frac{\sin x}{x\log^2x}\), the same as (a).
\end{enumerate}
\hint{3}
Answers (\(+C\)): (a) and (e) \(\frac{\sin x}{\log x}\); (b) \(\ln\lvert x+\log x\rvert\); (c) \(\ln\lvert x+\sin x\rvert\); (d) \(\ln\lvert x\rvert+\frac1{2022}\ln(1+x^{2022})\).
\end{document}
